\documentclass[11pt, letterpaper]{article}
\usepackage{mobi-lectura}

\title{Tiempos de Primera Pasada y Tiempos de Ocupación\\[4pt]
       {\large Cadenas de Markov de Tiempo Continuo}}
\author{Modelado de Sistemas bajo Incertidumbre}
\date{}

\begin{document}
\maketitle

\tableofcontents
\newpage

% ══════════════════════════════════════════════════════════════════════
\section{Preliminares: recordatorio de CTMC}
% ══════════════════════════════════════════════════════════════════════

Consideremos una \textbf{cadena de Markov de tiempo continuo} (CTMC)
$\{X(t),\; t\ge 0\}$ con espacio de estados finito $S=\{1,2,\ldots,N\}$.
Su dinámica queda completamente descrita por la \emph{matriz de tasas}
(o \emph{generador infinitesimal})
\[
  Q = [q_{ij}]_{i,j\in S},
\]
donde $q_{ij}\ge 0$ para $i\ne j$ y
\[
  q_{ii} = -\sum_{j\ne i} q_{ij} \eqqcolon -q_i, \qquad i\in S.
\]

Recordemos la interpretación: cuando la cadena se encuentra en el
estado~$i$, permanece allí un tiempo exponencial de tasa~$q_i$; al
abandonar~$i$, salta a~$j\ne i$ con probabilidad
\[
  p_{ij} = \frac{q_{ij}}{q_i}, \qquad j\ne i.
\]
La \emph{cadena embebida} (de tiempo discreto) tiene matriz de
transición $P=[p_{ij}]$ con $p_{ii}=0$.

\begin{notabox}[Notación]
  A lo largo de esta lectura, $q_{ij}$ denota la tasa de transición
  de~$i$ a~$j$; la tasa total de salida del estado~$i$ es
  $q_i = -q_{ii}$.  La matriz~$Q$ satisface $Q\,\mathbf{1}=\mathbf{0}$.
\end{notabox}

% ══════════════════════════════════════════════════════════════════════
\section{Tiempos de primera pasada}
% ══════════════════════════════════════════════════════════════════════

\subsection{Definición y motivación}

\begin{definicionbox}[Tiempo de primera pasada]
  Sea $\{X(t)\}$ una CTMC con espacio de estados~$S$ y sea
  $A\subseteq S$ un subconjunto no vacío.  Para cada estado
  inicial~$i\in S$ definimos el \textbf{tiempo de primera pasada}
  (o \textbf{tiempo de primer arribo}) al conjunto~$A$ como
  \[
    T_A = \inf\{t > 0 : X(t) \in A\}.
  \]
  En el caso particular $A=\{j\}$ escribimos simplemente $T_j$.
  Si $X(0)=j$ y $A=\{j\}$, entonces $T_j$ se interpreta como el
  \textbf{tiempo de retorno} al estado~$j$.
\end{definicionbox}

La cantidad central de interés es el \textbf{tiempo medio de primera
pasada} (\emph{mean first passage time}, MFPT):
\[
  m_{iA} \;=\; \E_i[T_A], \qquad i\in S,
\]
donde $\E_i[\cdot]$ denota la esperanza condicionada a $X(0)=i$.

\medskip
\textbf{¿Por qué es importante?}  Los MFPT permiten responder
preguntas como:
\begin{itemize}
  \item ¿Cuánto tiempo tarda, en promedio, un sistema en pasar de un
        modo de operación a otro?
  \item ¿Cuál es el tiempo medio de retorno a un estado específico?
  \item ¿Cuánto demora una cola en vaciarse por primera vez?
\end{itemize}

\subsection{Ecuaciones para el MFPT}

\begin{metodobox}[Teorema: Ecuaciones de primer paso para CTMC]
  Sea $A\subsetneq S$ y defínase $B = S\setminus A$.  Si la cadena
  es irreducible (o si $A$ es alcanzable desde todo estado en~$B$),
  entonces los tiempos medios de primera pasada $m_{iA}$, $i\in B$,
  satisfacen el sistema lineal
  \begin{equation}\label{eq:mfpt_sistema}
    \sum_{j\in B} q_{ij}\,m_{jA} = -1, \qquad i \in B,
  \end{equation}
  con la condición de frontera $m_{iA}=0$ para todo $i\in A$.

  \medskip
  En forma matricial, si $Q_B = [q_{ij}]_{i,j\in B}$ es la
  submatriz de~$Q$ restringida a~$B$:
  \begin{equation}\label{eq:mfpt_matricial}
    Q_B\,\mathbf{m} = -\mathbf{1}
    \qquad\Longrightarrow\qquad
    \mathbf{m} = -Q_B^{-1}\,\mathbf{1}.
  \end{equation}
\end{metodobox}

\textbf{Idea de la demostración (condicionamiento al primer paso).}
Dado $X(0)=i\in B$, la cadena permanece en~$i$ un tiempo exponencial
de media $1/q_i$.  Al abandonar~$i$ salta a algún estado~$j$ con
probabilidad $p_{ij}=q_{ij}/q_i$.  Por la propiedad de Markov
fuerte:
\[
  m_{iA} = \frac{1}{q_i} + \sum_{j\in B} p_{ij}\,m_{jA}
         = \frac{1}{q_i} + \sum_{j\in B} \frac{q_{ij}}{q_i}\,m_{jA}.
\]
Multiplicando ambos lados por $q_i$ y recordando que $q_{ii}=-q_i$
se obtiene~\eqref{eq:mfpt_sistema}. $\square$

\begin{notabox}[Invertibilidad de $Q_B$]
  Si $A$ es alcanzable desde todo estado de~$B$, la submatriz~$Q_B$
  es \emph{no singular}: todos sus valores propios tienen parte real
  estrictamente negativa.  Esto garantiza la existencia y unicidad
  de la solución~\eqref{eq:mfpt_matricial}.
\end{notabox}

\subsection{Ejemplo: MFPT en una cadena de 3 estados}

\begin{ejemplobox}[MFPT en una cadena de 3 estados]
  Consideremos una CTMC con $S=\{1,2,3\}$ y generador
  \[
    Q = \begin{pmatrix}
      -3 &  2 &  1 \\
       1 & -4 &  3 \\
       2 &  1 & -3
    \end{pmatrix}.
  \]

  \textbf{Pregunta:} ¿Cuánto tiempo tarda en promedio la cadena en
  llegar al estado~$3$ partiendo de cada uno de los otros estados?

  \medskip\noindent\textbf{Solución.}
  Aquí $A=\{3\}$ y $B=\{1,2\}$.  La submatriz $Q_B$ es
  \[
    Q_B = \begin{pmatrix} -3 & 2 \\ 1 & -4 \end{pmatrix}.
  \]
  Resolvemos $Q_B\,\mathbf{m}=-\mathbf{1}$:
  \[
    \det(Q_B) = (-3)(-4)-(2)(1) = 10, \quad
    Q_B^{-1} = \frac{1}{10}\begin{pmatrix} -4 & -2 \\ -1 & -3 \end{pmatrix}.
  \]
  Por tanto:
  \[
    \mathbf{m} = -Q_B^{-1}\mathbf{1}
    = -\frac{1}{10}\begin{pmatrix} -6 \\ -4 \end{pmatrix}
    = \begin{pmatrix} 3/5 \\ 2/5 \end{pmatrix}.
  \]
\end{ejemplobox}

\begin{resultadobox}[Resultado]
  $m_{13} = \dfrac{3}{5} = 0.6$ u.t. \quad y \quad $m_{23} = \dfrac{2}{5} = 0.4$ u.t.

  \medskip
  Es natural que $m_{23}<m_{13}$, pues desde el estado~2 la tasa
  directa hacia~3 es~3, mayor que la tasa~1 desde el estado~1.
\end{resultadobox}

% ══════════════════════════════════════════════════════════════════════
\section{Tiempos de ocupación}
% ══════════════════════════════════════════════════════════════════════

\subsection{Definición general y motivación}

\begin{definicionbox}[Tiempo de ocupación en horizonte fijo]
  Sea $\{X(t)\}$ una CTMC con espacio de estados~$S$, y sea~$j\in S$
  un estado fijo.  Para un horizonte de tiempo~$t>0$, el
  \textbf{tiempo de ocupación} del estado~$j$ en el intervalo
  $[0,t]$ es
  \[
    M_j(t) = \int_0^t \mathbf{1}_{\{X(s)=j\}}\,ds.
  \]
  Más generalmente, dado un conjunto $A\subseteq S$:
  \[
    M_A(t) = \int_0^t \mathbf{1}_{\{X(s)\in A\}}\,ds
            = \sum_{j\in A} M_j(t).
  \]
\end{definicionbox}

El tiempo de ocupación $M_j(t)$ mide \emph{cuánto tiempo acumulado}
pasa la cadena en el estado~$j$ durante $[0,t]$.  Es una variable
aleatoria con valores en~$[0,t]$ y, naturalmente,
\[
  \sum_{j\in S} M_j(t) = t \quad\text{con probabilidad 1.}
\]

\textbf{¿Por qué es importante?}
\begin{itemize}
  \item ¿Cuánto tiempo acumulado opera un servidor durante las
        primeras 10 horas?
  \item ¿Qué fracción de un periodo de observación pasa el sistema
        en estado de falla?
  \item ¿Cuál es el costo esperado acumulado si el costo depende
        del estado?
\end{itemize}

\subsection{Valor esperado del tiempo de ocupación: la fórmula integral}

\begin{metodobox}[Teorema: Tiempo medio de ocupación en horizonte fijo]
  Sea $p_{ij}(s) = \Prob_i(X(s)=j)$ la probabilidad de transición en
  tiempo~$s$.  Entonces
  \begin{equation}\label{eq:occ_esperanza}
    \E_i[M_j(t)] = \int_0^t p_{ij}(s)\,ds
    = \int_0^t \bigl[e^{Qs}\bigr]_{ij}\,ds.
  \end{equation}
  En forma matricial, la \textbf{matriz de tiempos medios de
  ocupación} es:
  \begin{equation}\label{eq:occ_matricial}
    \mathcal{O}(t) = \int_0^t e^{Qs}\,ds.
  \end{equation}
\end{metodobox}

\textbf{Demostración.}  Por linealidad de la esperanza y la
definición del indicador:
\[
  \E_i[M_j(t)] = \E_i\!\left[\int_0^t \mathbf{1}_{\{X(s)=j\}}\,ds\right]
  = \int_0^t \Prob_i(X(s)=j)\,ds
  = \int_0^t p_{ij}(s)\,ds.
\]
El intercambio esperanza--integral se justifica por Fubini.
Usando la solución de Kolmogorov hacia adelante, $P(s) = e^{Qs}$,
se obtiene~\eqref{eq:occ_esperanza}. $\square$

\begin{notabox}[Cómputo práctico vía diagonalización]
  Si $Q = V\Lambda V^{-1}$ con
  $\Lambda=\mathrm{diag}(\lambda_1,\ldots,\lambda_N)$, entonces:
  \[
    \mathcal{O}(t) = V\!\left[\mathrm{diag}\!\left(
    \frac{e^{\lambda_k t}-1}{\lambda_k}\right)\right]V^{-1},
  \]
  con la convención $\frac{e^{0\cdot t}-1}{0}=t$ para el valor propio
  $\lambda=0$ (que siempre existe en una CTMC irreducible).
\end{notabox}

\begin{notabox}[Comportamiento asintótico y teorema ergódico]
  Para $t$ grande, $e^{Qs}\to\mathbf{1}\,\boldsymbol{\pi}$ (cada fila
  converge a~$\boldsymbol{\pi}$, la distribución estacionaria), de modo que:
  \[
    \frac{\E_i[M_j(t)]}{t} \;\xrightarrow{t\to\infty}\; \pi_j,
  \]
  independientemente del estado inicial.  Esta es la manifestación
  del \textbf{teorema ergódico}: la fracción de tiempo que la cadena
  pasa en cada estado converge a la probabilidad estacionaria.
\end{notabox}

\subsection{Ejemplo: tiempos de ocupación en cadena de 2 estados}

\begin{ejemplobox}[Tiempos medios de ocupación para una cadena de 2 estados]
  Consideremos $S=\{1,2\}$ y generador
  \[
    Q = \begin{pmatrix} -\alpha & \alpha \\ \beta & -\beta \end{pmatrix},
    \qquad \alpha,\beta>0.
  \]
  Los valores propios son $\lambda_1=0$ y $\lambda_2=-(\alpha+\beta)$.
  Integrando componente a componente en $[0,t]$:
  \begin{align*}
    \E_1[M_1(t)] &= \frac{1}{\alpha+\beta}\left[\beta\,t
    + \frac{\alpha}{\alpha+\beta}\bigl(1-e^{-(\alpha+\beta)t}\bigr)\right],\\[4pt]
    \E_1[M_2(t)] &= \frac{1}{\alpha+\beta}\left[\alpha\,t
    - \frac{\alpha}{\alpha+\beta}\bigl(1-e^{-(\alpha+\beta)t}\bigr)\right].
  \end{align*}

  \textbf{Verificación:} $\E_1[M_1(t)]+\E_1[M_2(t)] = t\;\checkmark$

  \medskip
  \textbf{Caso numérico:} $\alpha=2$, $\beta=3$, $t=1$.
  La distribución estacionaria es $\pi_1=3/5=0.6$,
  $\pi_2=2/5=0.4$.
  \[
    \E_1[M_1(1)] \approx 0.680, \qquad \E_1[M_2(1)] = 0.320.
  \]
  Compare con $\pi_1\cdot t = 0.6$: en $t=1$ la cadena aún no
  ha alcanzado el equilibrio (partiendo del estado~1, todavía pasa
  más tiempo en~1 que lo que dictaría~$\pi_1$).
\end{ejemplobox}

\subsection{Tiempos de ocupación hasta la primera pasada}

Hasta ahora hemos estudiado los tiempos de ocupación en un horizonte
fijo~$t$.  Un caso particular importante surge cuando el horizonte es
el tiempo de primera pasada $T_A$ a un conjunto~$A$.

\begin{definicionbox}[Tiempo de ocupación hasta primera pasada]
  Sea $A\subsetneq S$ y $B=S\setminus A$.  El \textbf{tiempo de
  ocupación del estado~$j\in B$ hasta la primera pasada a~$A$} es
  \[
    M_j^{(A)} = \int_0^{T_A} \mathbf{1}_{\{X(s)=j\}}\,ds,
    \qquad j\in B.
  \]
\end{definicionbox}

\begin{metodobox}[Teorema: Ecuaciones para tiempos medios de ocupación hasta primera pasada]
  Definamos $\mu_{ij} = \E_i[M_j^{(A)}]$ para $i,j\in B$.  Entonces
  los valores $\{\mu_{ij}\}_{i\in B}$ satisfacen, para~$j$ fijo:
  \begin{equation}\label{eq:occ_fp}
    \sum_{k\in B} q_{ik}\,\mu_{kj} = -\delta_{ij},
    \qquad i\in B.
  \end{equation}
  En forma matricial:
  \begin{equation}\label{eq:occ_fp_mat}
    Q_B\,\mathcal{M} = -I_{|B|}
    \qquad\Longrightarrow\qquad
    \mathcal{M} = -Q_B^{-1},
  \end{equation}
  donde $\mathcal{M}=[\mu_{ij}]_{i,j\in B}$.
\end{metodobox}

\begin{notabox}[Relación entre los dos enfoques]
  \begin{itemize}
    \item \textbf{Horizonte fijo~$t$}: $\mathcal{O}(t)=\int_0^t e^{Qs}\,ds$
      — requiere la exponencial matricial; el resultado es función de~$t$.
    \item \textbf{Hasta primera pasada}: $\mathcal{M}=-Q_B^{-1}$
      — solo álgebra lineal; el resultado es un número fijo.
  \end{itemize}
  El segundo caso es algebraicamente más simple, pero es una
  especialización del primero con horizonte aleatorio~$T_A$.
\end{notabox}

\begin{notabox}[Conexión MFPT $\leftrightarrow$ ocupación hasta primera pasada]
  El MFPT se expresa como la suma de todos los tiempos medios de
  ocupación en los estados de~$B$:
  \[
    m_{iA} = \E_i[T_A] = \sum_{j\in B} \mu_{ij}
    \qquad\Longleftrightarrow\qquad
    \mathbf{m} = \mathcal{M}\,\mathbf{1} = (-Q_B^{-1})\,\mathbf{1}.
  \]
  La matriz $-Q_B^{-1}$ contiene \emph{simultáneamente} toda la
  información sobre los tiempos de ocupación y los MFPT.
\end{notabox}

\begin{ejemplobox}[Tiempos de ocupación hasta primera pasada en la cadena de 3 estados]
  Con la misma cadena ($S=\{1,2,3\}$, $A=\{3\}$, $B=\{1,2\}$):
  \[
    \mathcal{M} = -Q_B^{-1}
    = \frac{1}{10}\begin{pmatrix} 4 & 2 \\ 1 & 3 \end{pmatrix}.
  \]

  \textbf{Interpretación:}
  \begin{itemize}
    \item $\mu_{11}=4/10$: partiendo de~1, tiempo medio en el estado~1 antes de llegar a~3.
    \item $\mu_{12}=2/10$: partiendo de~1, tiempo medio en el estado~2 antes de llegar a~3.
    \item $\mu_{21}=1/10$: partiendo de~2, tiempo medio en el estado~1 antes de llegar a~3.
    \item $\mu_{22}=3/10$: partiendo de~2, tiempo medio en el estado~2 antes de llegar a~3.
  \end{itemize}

  \textbf{Verificación con el MFPT:}
  \[
    m_{13} = \mu_{11}+\mu_{12} = \tfrac{4}{10}+\tfrac{2}{10}
    = \tfrac{3}{5}\;\checkmark,
    \qquad
    m_{23} = \mu_{21}+\mu_{22} = \tfrac{1}{10}+\tfrac{3}{10}
    = \tfrac{2}{5}\;\checkmark.
  \]
\end{ejemplobox}

% ══════════════════════════════════════════════════════════════════════
\section{Ejemplo detallado con dos variables de estado}
% ══════════════════════════════════════════════════════════════════════

Consideremos un sistema con espacio de estados $S=\{0,1,2\}$ descrito
por dos variables indicadoras: $X$ indica si el componente está
operativo ($X=1$) o no ($X=0$), e $Y$ indica si hay respaldo
disponible ($Y=1$) o no ($Y=0$).

\begin{center}
\begin{tabular}{ccl}
  \toprule
  \textbf{Estado} & $(X,Y)$ & \textbf{Descripción}\\
  \midrule
  0 & $(0,0)$ & Componente inactivo, sin respaldo\\
  1 & $(1,0)$ & Componente operativo, sin respaldo\\
  2 & $(1,1)$ & Componente operativo, con respaldo\\
  \bottomrule
\end{tabular}
\end{center}

\subsection{Generador infinitesimal}

\[
  Q = \begin{pmatrix}
    -2 &  2 &  0 \\
     1 & -5 &  4 \\
     0 &  3 & -3
  \end{pmatrix}.
\]

Nótese que el estado~$0$ no se comunica directamente con el
estado~$2$: para llegar al modo completo hay que pasar primero por el
modo operativo sin respaldo.

\subsection{Cálculo de los tiempos medios de primera pasada}

\textbf{Pregunta 1: ¿Cuánto tarda en promedio el sistema en alcanzar
el modo completo (estado~2) desde cada estado?}

Con $A=\{2\}$, $B=\{0,1\}$, $Q_B = \begin{pmatrix} -2 & 2 \\ 1 & -5 \end{pmatrix}$,
$\det(Q_B)=8$:
\[
  \mathbf{m} = -Q_B^{-1}\,\mathbf{1}
  = \begin{pmatrix} 7/8 \\ 3/8 \end{pmatrix}.
\]

\begin{resultadobox}[MFPT al estado 2]
  $m_{02} = \dfrac{7}{8} = 0.875$ u.t., \quad
  $m_{12} = \dfrac{3}{8} = 0.375$ u.t.
\end{resultadobox}

\medskip
\textbf{Pregunta 2: ¿Cuánto tarda en promedio en llegar al estado
inactivo (estado~0) partiendo del modo completo?}

Con $A=\{0\}$, $B=\{1,2\}$, $Q_B = \begin{pmatrix} -5 & 4 \\ 3 & -3 \end{pmatrix}$,
$\det(Q_B)=3$:
\[
  \mathbf{m} = -Q_B^{-1}\,\mathbf{1}
  = \begin{pmatrix} 7/3 \\ 8/3 \end{pmatrix}.
\]

\begin{resultadobox}[MFPT al estado 0]
  $m_{10} = \dfrac{7}{3} \approx 2.333$ u.t., \quad
  $m_{20} = \dfrac{8}{3} \approx 2.667$ u.t.
\end{resultadobox}

Es razonable que el estado~2 (modo completo) demore más en caer al
estado~0 que el estado~1, pues primero debe perder el respaldo y
luego fallar.

\subsection{Cálculo de los tiempos medios de ocupación}

\textbf{Pregunta 3: Antes de alcanzar el modo completo (estado~2),
¿cuánto tiempo pasa el sistema en cada estado?}

Con $A=\{2\}$, $B=\{0,1\}$:
\[
  \mathcal{M} = -Q_B^{-1}
  = \frac{1}{8}\begin{pmatrix} 5 & 2 \\ 1 & 2 \end{pmatrix}.
\]

\begin{resultadobox}[Tiempos medios de ocupación antes de llegar al estado 2]
\begin{center}
\renewcommand{\arraystretch}{1.4}
\begin{tabular}{c|cc|c}
  \toprule
  \textbf{Desde} & $\mu_{i0}$ (en estado~0) & $\mu_{i1}$ (en estado~1)
  & $m_{i2}=\mu_{i0}+\mu_{i1}$\\
  \midrule
  $i=0$ & $5/8=0.625$ & $2/8=0.250$ & $7/8=0.875$\\
  $i=1$ & $1/8=0.125$ & $2/8=0.250$ & $3/8=0.375$\\
  \bottomrule
\end{tabular}
\end{center}
\end{resultadobox}

\medskip
\textbf{Pregunta 4: Antes de caer al estado inactivo (estado~0),
¿cuánto tiempo pasa en cada modo operativo?}

Con $A=\{0\}$, $B=\{1,2\}$:
\[
  \mathcal{M} = -Q_B^{-1}
  = \frac{1}{3}\begin{pmatrix} 3 & 4 \\ 3 & 5 \end{pmatrix}
  = \begin{pmatrix} 1 & 4/3 \\ 1 & 5/3 \end{pmatrix}.
\]

\begin{resultadobox}[Tiempos medios de ocupación antes de llegar al estado 0]
\begin{center}
\renewcommand{\arraystretch}{1.4}
\begin{tabular}{c|cc|c}
  \toprule
  \textbf{Desde} & $\mu_{i1}$ (operativo) & $\mu_{i2}$ (completo)
  & $m_{i0}=\mu_{i1}+\mu_{i2}$\\
  \midrule
  $i=1$ & $1$ & $4/3\approx 1.333$ & $7/3\approx 2.333$\\
  $i=2$ & $1$ & $5/3\approx 1.667$ & $8/3\approx 2.667$\\
  \bottomrule
\end{tabular}
\end{center}
\end{resultadobox}

\begin{notabox}[Observación notable]
  Sin importar si se parte del estado~1 o del estado~2, el tiempo
  medio de ocupación en el estado~1 antes de caer al estado~0 es el
  mismo ($\mu_{11}=\mu_{21}=1$).  Esto se explica porque desde el
  estado~2 la cadena \emph{necesariamente} pasa por~1 antes de llegar
  a~0 y, al llegar a~1, la propiedad de Markov ``reinicia'' el proceso.
\end{notabox}

% ══════════════════════════════════════════════════════════════════════
\section{Preguntas de aplicación con valores esperados}
% ══════════════════════════════════════════════════════════════════════

Las siguientes preguntas se refieren al sistema de la sección anterior.

\begin{ejemplobox}[Ejercicio: Valor esperado del costo hasta la falla total]
  El sistema incurre en un costo por unidad de tiempo de \$5 cuando
  está en el estado~1 (operativo sin respaldo) y de \$2 cuando está
  en el estado~2 (completo).  Sea~$C$ el costo total acumulado desde
  que el sistema arranca en modo completo (estado~2) hasta que cae al
  estado~0 (falla total).  Calcule $\E[C]$.

  \medskip
  \textbf{Solución.}  El costo total es
  $C = 5\,M_1^{(\{0\})} + 2\,M_2^{(\{0\})}$.
  Por linealidad de la esperanza, con estado inicial~2:
  \[
    \E_2[C] = 5\,\mu_{21} + 2\,\mu_{22}
    = 5\cdot 1 + 2\cdot\frac{5}{3}
    = \frac{25}{3} \approx \$8.33.
  \]
\end{ejemplobox}

\begin{ejemplobox}[Ejercicio: Número esperado de visitas a un estado]
  ¿Cuántas veces se espera que el sistema visite el estado~1 antes de
  alcanzar el modo completo (estado~2), si arranca en el estado~0?

  \medskip
  \textbf{Solución.}  El número esperado de visitas al estado~$j$ está
  relacionado con el tiempo de ocupación.  Cada visita al estado~$j$
  dura un tiempo exponencial de media $1/q_j$:
  \[
    \E_i[M_j^{(A)}] = \E_i[N_j]\cdot\frac{1}{q_j}
    \qquad\Longrightarrow\qquad
    \E_i[N_j] = q_j\,\mu_{ij}.
  \]
  Con $i=0$, $j=1$, $A=\{2\}$: $q_1=5$ y $\mu_{01}=2/8=1/4$:
  \[
    \E_0[N_1] = 5\cdot\frac{1}{4} = \frac{5}{4} = 1.25 \text{ visitas.}
  \]
\end{ejemplobox}

\begin{ejemplobox}[Ejercicio: Tiempo medio de retorno en régimen estacionario]
  Encuentre la distribución estacionaria de la cadena y verifique
  que el tiempo medio de retorno al estado~$j$ satisface
  $\E_j[T_j] = 1/(\pi_j\,q_j)$.

  \medskip
  \textbf{Solución.}  Resolviendo $\boldsymbol{\pi}\,Q=\mathbf{0}$
  con $\sum\pi_j=1$:
  \[
    \pi_0 = \frac{3}{17}\approx 0.176, \quad
    \pi_1 = \frac{6}{17}\approx 0.353, \quad
    \pi_2 = \frac{8}{17}\approx 0.471.
  \]
  Para el estado~0: $\E_0[T_0] = 1/(\pi_0\,q_0) = 1/((3/17)(2)) = 17/6 \approx 2.833$ u.t.

  Para el estado~2: $\E_2[T_2] = 1/(\pi_2\,q_2) = 1/((8/17)(3)) = 17/24 \approx 0.708$ u.t.
\end{ejemplobox}

% ══════════════════════════════════════════════════════════════════════
\section{Resumen de fórmulas clave}
% ══════════════════════════════════════════════════════════════════════

Sea $\{X(t)\}$ una CTMC irreducible con generador~$Q$ sobre
$S=\{1,\ldots,N\}$.  Sea $A\subsetneq S$, $B=S\setminus A$,
y $Q_B$ la submatriz de~$Q$ con filas y columnas en~$B$.

\begin{metodobox}[Fórmulas principales]

\textbf{1.\ Tiempos medios de primera pasada} ($m_{iA}=\E_i[T_A]$,
  $i\in B$):
  \[
    Q_B\,\mathbf{m} = -\mathbf{1}
    \qquad\Longrightarrow\qquad
    \mathbf{m} = -Q_B^{-1}\,\mathbf{1}.
  \]

\textbf{2.\ Tiempos medios de ocupación en horizonte fijo~$t$}:
  \[
    \E_i[M_j(t)] = \int_0^t [e^{Qs}]_{ij}\,ds,
    \qquad
    \mathcal{O}(t) = \int_0^t e^{Qs}\,ds.
  \]

\textbf{3.\ Tiempos medios de ocupación hasta primera pasada}
  ($\mu_{ij}=\E_i[M_j^{(A)}]$, $i,j\in B$):
  \[
    Q_B\,\mathcal{M} = -I
    \qquad\Longrightarrow\qquad
    \mathcal{M} = -Q_B^{-1}.
  \]

\textbf{4.\ Relación MFPT $\leftrightarrow$ ocupación hasta primera
  pasada:}
  \[
    m_{iA} = \sum_{j\in B}\mu_{ij}, \qquad
    \mathbf{m} = \mathcal{M}\,\mathbf{1}.
  \]

\textbf{5.\ Número esperado de visitas} al estado~$j\in B$ antes de
  la primera pasada a~$A$, partiendo de~$i\in B$:
  \[
    \E_i[N_j] = q_j\,\mu_{ij}.
  \]

\textbf{6.\ Comportamiento asintótico} ($t\to\infty$, cadena
  irreducible):
  \[
    \frac{\E_i[M_j(t)]}{t}\;\xrightarrow{t\to\infty}\;\pi_j.
  \]

\textbf{7.\ Tiempo medio de retorno} al estado~$j$:
  \[
    \E_j[T_j] = \frac{1}{\pi_j\,q_j}.
  \]

\end{metodobox}

% ══════════════════════════════════════════════════════════════════════
\section*{Referencias}
% ══════════════════════════════════════════════════════════════════════

\begin{enumerate}
  \item V.\,G.\ Kulkarni, \textit{Modeling and Analysis of Stochastic
    Systems}, 3.ª ed., CRC Press, 2017.  Capítulos~6 y~8.
  \item J.\,R.\ Norris, \textit{Markov Chains}, Cambridge University
    Press, 1997.
  \item W.\,J.\ Stewart, \textit{Probability, Markov Chains, Queues,
    and Simulation}, Princeton University Press, 2009.
\end{enumerate}

\end{document}
